% Complete independently derived fragment.  Requires amsmath, amssymb,
% amsthm, amscd, hyperref, and theorem/proposition/lemma environments.
% Tags HG1--HG41 and HG21a--HG21c are reserved for this fragment.
\section{The full Split-Zero coefficient object as a Connes--Consani
\texorpdfstring{$\Gamma$}{Gamma}-algebra}

The construction used here is Alain Connes and Caterina Consani,
\emph{Absolute algebra and Segal's $\Gamma$-rings: au dessous de
$\overline{\operatorname{Spec}\mathbb Z}$}.

The original author source is
\nolinkurl{F1_corr.tex}, associated with
\href{https://arxiv.org/abs/1502.05585v2}{arXiv:1502.05585v2}.
The whole source, including its bibliography and final empty lines,
was read: lines 1--1336.  The read file has SHA-256
\begin{center}\small\texttt{3b1319ec1df4ed6e2412ace960c2813b683a}\\
\texttt{36c88509234c9abd0a39f98e8913}.\end{center}
The local source has duplicate labels \texttt{sssalg1}; below the
line numbers distinguish the full-faithfulness proposition (line 536)
from the Boolean-semiring lemma (line 575).  The formula
\texttt{functsum}, lines 433--442, and Lemma \texttt{sssalg},
lines 519--533, specify the actual functor and multiplication.
Proposition \texttt{sssalg2}, lines 905--932, specifies the norm
construction.  The conclusions below are direct calculations with
those formulas, not an identification of every spherical algebra with
a semiring.

Let $R$ be a nonzero commutative ring with identity and let $L$ be a
bounded distributive lattice.  Write its bottom and top as $0_L,1_L$.
Keep the entire coefficient object
\begin{equation}\tag{HG1}
 \begin{gathered}
 G=G_L(R)=\{(0,\lambda):\lambda\in L\}
       \cup\{(r,1_L):r\in R\},\\
 (r,\lambda)+(s,\mu)=(r+s,\lambda\vee\mu),\qquad
 (r,\lambda)(s,\mu)=(rs,\lambda\wedge\mu),\\
 \tau=(0,0_L),\quad e=(0,1_L),\quad
 \widehat r=(r,1_L),\quad z_\lambda=(0,\lambda).
 \end{gathered}
\end{equation}
Thus $\widehat0=e$.  No occurrence of $e$ in this fragment is
silently replaced by $\tau$.  The semiring identity is $\widehat1$.
The formulas close on the stated carrier, and associativity and
distributivity follow coordinate by coordinate from the ring and
lattice laws.  Regard $L$ itself as the semiring with addition $\vee$
and multiplication $\wedge$.

\subsection{Every finite pointed set, every fold, and multiplication}

For a finite pointed set $X$ let $X^\circ=X\setminus\{*\}$.
The base point of $HG(X)$ is the tuple constantly equal to $\tau$.
For $X=0_+$ it is the unique empty tuple.  The exact formula is
\begin{equation}\tag{HG2}
 HG(X)=\left\{((r_x,\lambda_x))_{x\in X^\circ}:
       r_x\ne0\Longrightarrow\lambda_x=1_L\right\}.
\end{equation}
For a pointed map $f:X\to Y$ set
\begin{equation}\tag{HG3}
 (HG(f)(r,\lambda))_y=
 \left(\sum_{x\in X^\circ:f(x)=y}r_x,
       \bigvee_{x\in X^\circ:f(x)=y}\lambda_x\right),
 \qquad y\in Y^\circ.
\end{equation}
An empty sum is $0_R$ and an empty join is $0_L$.  Coordinates sent
to the base point of $Y$ are discarded, rather than added to a
retained coordinate.

\begin{theorem}\label{HG:functor}
Equations (HG2)--(HG3), together with
\begin{equation}\tag{HG4}
 \mu_{X,Y}((r,\lambda),(s,\mu))_{(x,y)}
          =(r_xs_y,\lambda_x\wedge\mu_y),
 \quad (x,y)\in X^\circ\times Y^\circ,
\end{equation}
define the Connes--Consani commutative $\mathbb S$-algebra $HG$.
Its unit $\mathbb S(X)=X\to HG(X)$ sends $x\ne *$ to the
tuple with $\widehat1$ at $x$ and $\tau$ elsewhere, and sends
$*$ to the all-$\tau$ tuple.
\end{theorem}
\begin{proof}
If a sum in (HG3) is nonzero, at least one summand is nonzero,
so its support is $1_L$ and the output join is $1_L$.  Thus the
output satisfies (HG2), even when other terms cancel.  Identity
maps act identically.  For $g:Y\to Z$ and $z\ne *$, the fibres
$f^{-1}(y)$ for $g(y)=z$ partition $(gf)^{-1}(z)$; no point
discarded by $f$ can reappear, since $g(*)=*$.  Associativity
of finite sums and joins now gives $HG(g)HG(f)=HG(gf)$.

Formula (HG4) lands in (HG2) because a nonzero product requires
both amplitudes to be nonzero.  An all-$\tau$ input produces an
all-$\tau$ output, so this map factors through the smash product
of pointed sets.  Naturality in $f:X\to X'$ and $g:Y\to Y'$
is the equality, at $(x',y')$, between the sum over
$f^{-1}(x')\times g^{-1}(y')$ of the products and the product
of the two fibre sums.  The amplitude equality is ring
distributivity; the support equality is
\[
 \bigvee_{x,y}(\lambda_x\wedge\mu_y)
       =\left(\bigvee_x\lambda_x\right)
          \wedge\left(\bigvee_y\mu_y\right).
\]
This natural bifunctor map is precisely the multiplication map
for the Day smash product, by the source's definition preceding
\texttt{colim}, lines 445--464.  The two triple products agree
coordinatewise, as do the switched products.  The unit rule is
natural because each unit tuple has at most one non-$\tau$
coordinate; the product with that tuple has its original value
on its single selected coordinate and $\tau$ elsewhere.  This
proves all claimed algebra axioms directly.
\end{proof}

\subsection{Amplitude and lattice support form one exact object}

Define $q:G\to R$ and $p:G\to L$ by
\begin{equation}\tag{HG5}
 q(r,\lambda)=r,\qquad p(r,\lambda)=\lambda.
\end{equation}
Both maps preserve zero, identity, addition and multiplication.
Consequently $Hq:HG\to HR$ and $Hp:HG\to HL$ are morphisms
of $\mathbb S$-algebras, not merely maps on level $1$.

\begin{theorem}\label{HG:joint}
The joint map
\begin{equation}\tag{HG6}
 (Hq,Hp):HG\hookrightarrow HR\times HL
\end{equation}
is injective at every finite pointed set.  Its image is exactly
\begin{equation}\tag{HG7}
 \{(r,\lambda)\in R^{X^\circ}\times L^{X^\circ}:
             r_x\ne0\Longrightarrow\lambda_x=1_L\}.
\end{equation}
For fixed amplitude $r$ and fixed support $\lambda$, respectively,
the fibres are
\begin{equation}\tag{HG8}
 \begin{aligned}
 (Hq)_X^{-1}(r)&\cong\prod_{x:r_x=0}L,\\
 (Hp)_X^{-1}(\lambda)&\cong
                     \prod_{x:\lambda_x=1_L}R.
 \end{aligned}
\end{equation}
All omitted support coordinates in the first line are $1_L$;
all omitted amplitude coordinates in the second line are $0_R$.
These descriptions also cover $0_L=1_L$, when $G_L(R)=R$.
\end{theorem}
\begin{proof}
A coefficient is exactly its amplitude and support pair, proving
injectivity and the image claim.  The naturality of the two
projections follows from (HG3): the first takes fibre sums and
the second takes fibre joins, which are the defining actions of
$HR$ and $HL$.  Equation (HG4) proves multiplication compatibility
and the unit tuples give unit compatibility.  At each zero
amplitude coordinate every support is allowed; at each nonzero
amplitude coordinate only top is allowed.  Conversely, at a
support different from top the amplitude must be zero, and at
top every amplitude is allowed.  Independence of the coordinates
proves (HG8).
\end{proof}

The pointwise inverse image of zero under $Hq$ is the embedded
full lattice object
\begin{equation}\tag{HG9}
 \iota:HL\longrightarrow HG,\quad
       (\lambda_x)_x\longmapsto(z_{\lambda_x})_x,
 \qquad Hp\circ\iota=\operatorname{id}_{HL}.
\end{equation}
This is a natural injective map of $\Gamma$-sets and preserves
products, but when $R\ne0$ it is not a unital map of
$\mathbb S$-algebras: its unit is sent to $e$ rather than
$\widehat1$.  Its image is an ideal for (HG4), with the explicit
action
\begin{equation}\tag{HG10}
 (r_x,\lambda_x)z_{\mu_y}=z_{\lambda_x\wedge\mu_y}.
\end{equation}
The map $Hp$ has only the base point in its zero fibre when
$0_L\ne1_L$, yet it is not injective: $e$ and $\widehat1$ have
the same support.  No argument appropriate only to abelian-group
exact sequences can erase these fibres.

\begin{proposition}\label{HG:fibretransport}
Fix $r\in HR(X)$ and $f:X\to Y$.  Write $r'=HR(f)r$ and
$I_y=\{x\in X^\circ:f(x)=y\}$.  The image of the entire
fibre $(Hq)_X^{-1}(r)$ in $(Hq)_Y^{-1}(r')$ consists exactly
of the following support choices, independently at each $y$:
\begin{equation}\tag{HG11}
 \lambda'_y\in
 \begin{cases}
 \{0_L\},& I_y=\varnothing,\\
 \{1_L\},&\text{some }x\in I_y\text{ has }r_x\ne0,\\
 L,&I_y\ne\varnothing\text{ and all }r_x=0\text{ on }I_y.
 \end{cases}
\end{equation}
In particular nontrivial cancellation in a fold produces $e$,
not $\tau$, while a fold of labelled zeros produces their exact
lattice join.
\end{proposition}
\begin{proof}
The first case is the empty-join convention.  In the second case
at least one input support is top, so the output support is top
regardless of whether $r'_y$ vanishes.  In the third case all
input supports are unrestricted.  Every $\lambda\in L$ occurs
as the join by assigning it to any one input coordinate and
assigning bottom to the others.  Distinct $I_y$ are disjoint,
so all these choices can be made independently.  The last
assertions are exactly (HG3).
\end{proof}

\subsection{The failed section and its complete correction object}

There is always the pointed set section
\begin{equation}\tag{HG12}
 b:R\to G,\qquad b(0)=\tau,\quad b(r)=\widehat r\ (r\ne0),
 \qquad q\circ b=\operatorname{id}_R.
\end{equation}
Let $b_X:R^{X^\circ}\to G^{X^\circ}$ apply this formula
coordinatewise.  This is a family of pointed maps, not yet a
natural transformation.

\begin{theorem}\label{HG:nosection}
Assume $0_L\ne1_L$.  Every natural transformation of
$\Gamma$-sets $HR\to HG$ is the constant base-point map.
In particular (HG12) does not lift to a $\Gamma$-map, and
$Hq$ has no $\Gamma$-section.  The pointed section $b$ is
multiplicative precisely when $R$ has no nonzero zero divisors.
\end{theorem}
\begin{proof}
Let $\eta:HR\to HG$ be natural and write $u=\eta_{1_+}$.
For each $x\in X^\circ$, the pointed projection
$p_x:X\to1_+$ takes $x$ to $1$ and every other point to $0$.
Naturality forces $(\eta_X(r))_x=u(r_x)$.  Apply this first to
$X=2_+$ and then to its fold $1,2\mapsto1$.  One obtains
$u(a+b)=u(a)+u(b)$, with $u(0)=\tau$.  This is the same
recovery of addition used by Connes--Consani in (\texttt{addition})
and the proof at lines 536--558, here proved for a $\Gamma$-map
without assuming it preserves multiplication.

If $v+w=\tau$ in $G$, their supports satisfy
$p(v)\vee p(w)=0_L$, so both supports are bottom.  Since
$0_L\ne1_L$, the carrier (HG1) then forces $v=w=\tau$.
Thus $\tau$ is the only additively invertible element of $G$.
For every $r\in R$ the equality $u(r)+u(-r)=u(0)=\tau$
forces $u(r)=\tau$.  The coordinate projections now force
every $\eta_X$ to be constant.

For multiplicativity of $b$, if either input vanishes the
identity follows from absorption by $\tau$.  If both are
nonzero and their product is nonzero, both sides equal the
top-supported product.  If $r,s\ne0$ but $rs=0$, the left
side $b(r)b(s)$ is $e$ and the right side $b(rs)$ is $\tau$.
This proves exactly the asserted criterion.
\end{proof}

For the obstruction formulas (HG13)--(HG16), take $0_L\ne1_L$.
When $L$ has one element, $G_L(R)=R$, $b$ is the identity,
and each correction is the unique zero element of $HL$.
To calculate the obstruction at every arity, use the two-element
subset $\{0_L,1_L\}$ inside $L$.  For $a\in R^{X^\circ}$
put $\nu(a)_x=1_L$ when $a_x\ne0$ and $0_L$ otherwise.
For a pointed map $f:X\to Y$ define
\begin{equation}\tag{HG13}
 \begin{aligned}
 C_f(a)_y=1_L\quad\Longleftrightarrow\quad
 &\sum_{x:f(x)=y}a_x=0\ \text{and}\\
 &\text{some }x\in X^\circ\text{ with }f(x)=y\text{ has }a_x\ne0.
 \end{aligned}
\end{equation}
otherwise put $C_f(a)_y=0_L$.  This defines an actual element of
$HL(Y)$, including empty fibres and discarded coordinates.
The exact failure-of-naturality equation is
\begin{equation}\tag{HG14}
 HG(f)b_X(a)=b_Y(HR(f)a)+\iota_Y(C_f(a)).
\end{equation}
Indeed both amplitudes are the fibre sum.  If that sum is nonzero,
both supports are top.  If it is zero, the left support is top
exactly when at least one input amplitude is nonzero, which is
precisely (HG13).  In particular, for $a\ne0$ the binary fold
of $(b(a),b(-a))$ is $e$ and its correction is $z_{1_L}$.

\begin{figure}[ht]
\centering
\[
\begin{aligned}
 (a,-a)&\xrightarrow{\ b_{2_+}\ }
      (\widehat a,\widehat{-a})
      \xrightarrow{\ HG(\mathrm{fold})\ }e=(0,1_L),\\[4pt]
 (a,-a)&\xrightarrow{\ HR(\mathrm{fold})\ }0
      \xrightarrow{\ b_{1_+}\ }\tau=(0,0_L).
\end{aligned}
\]
\caption{The two exact paths compared by naturality, for any
$a\ne0$ and $0_L\ne1_L$.  Their endpoints are different.
Equation (HG14) supplies their full correction inside the
embedded $HL$, and (HG15) computes its behavior under further
folds.  The fold operation is Connes--Consani's
\texttt{functsum}, source lines 433--442.  This is a diagram
of the calculated noncommutation, not a commuting square.}
\end{figure}

The correction itself has an exact composition rule.  Put
$Z(a)_x=1_L$ when $a_x=0$, and bottom otherwise.  For
$g:Y\to Z$ one has
\begin{equation}\tag{HG15}
 C_{gf}(a)=C_g(HR(f)a)\ \vee\
 \bigl(Z(HR(gf)a)\wedge HL(g)(C_f(a))\bigr),
\end{equation}
where meet and join are coordinatewise.  To prove it, a final
coordinate with nonzero amplitude has zero on both sides.  At
a final coordinate with zero amplitude, an original nonzero
summand is either in an intermediate fibre whose amplitude is
nonzero, or in an intermediate fibre that cancelled to zero.
The first event is exactly the first term on the right, and
the second is exactly its second term.  This includes their
possible simultaneous occurrence and proves (HG15).

There is also a complete multiplication correction for general
rings with zero divisors:
\begin{equation}\tag{HG16}
 \begin{gathered}
 D(a,b)_{x,y}=1_L\quad\Longleftrightarrow\quad
 a_x\ne0,\ b_y\ne0,\ a_xb_y=0,\\
 \mu(b_X(a),b_Y(b))=
 b_{X\wedge Y}((a_xb_y)_{x,y})+\iota_{X\wedge Y}(D(a,b)),
 \end{gathered}
\end{equation}
with $D=0_L$ at all other coordinates.  At a nonzero product
both sides have top support; at a zero product only the case
where both factors were nonzero contributes $e$.  This proves
the formula without a domain assumption.

\subsection{The exact spherical replacement and universal ring image}

Let $M(R)$ be the multiplicative monoid of $R$, pointed by its
ring zero.  The source's spherical monoid algebra is
$\mathbb S M(R)(X)=M(R)\wedge X$; see
\texttt{mono2sss} and \texttt{monomono}, lines 489--511.
The pointed family
\begin{equation}\tag{HG17}
 \widetilde b_X:M(R)\wedge X\to HG(X),\qquad
 (r,x)\longmapsto\bigl[\text{$b(r)$ at $x$, $\tau$ elsewhere}\bigr]
\end{equation}
is a $\Gamma$-map for every $R$: an ordinary pointed map either
retains the single selected coordinate, carrying its value with
it, or discards it entirely.  Its compatibility with the
multiplication of spherical algebras is equivalent to
$b(rs)=b(r)b(s)$ at level $1$, hence holds exactly under the
criterion in Theorem~\ref{HG:nosection}.  It is a unital
$\mathbb S$-algebra map for a domain.  When $0_L\ne1_L$
it is not such a map for a ring with nonzero zero divisors;
when $L$ has one element it is the source's usual map
$\mathbb S M(R)\to HR$ for every ring.

For every ring, including rings with zero divisors, the required
multiplicative replacement retains the two zero states.  Let
\begin{equation}\tag{HG18}
 M^+(R)=R\sqcup\{\perp\},\qquad
 \perp m=m\perp=\perp,
\end{equation}
with the original ring product on $R$, identity $1_R$, and
distinguished pointed zero $\perp$.  The original element
$0_R\in R$ is not $\perp$.  Define
\begin{equation}\tag{HG19}
 c:M^+(R)\to G,\quad c(\perp)=\tau,\quad c(r)=\widehat r,
 \qquad d:M^+(R)\to M(R),\quad d(\perp)=0,\quad d(r)=r.
\end{equation}
Both are pointed unital multiplicative maps, including products
of two nonzero zero divisors, because such a product lands at
the original $0_R$ and hence at $e$ under $c$.
The source adjunction gives the actual commuting diagram
\begin{equation}\tag{HG20}
 \begin{CD}
 \mathbb S M^+(R) @>{\widetilde c}>> HG\\
 @V{\mathbb S d}VV @VV{Hq}V\\
 \mathbb S M(R) @>{\rho_R}>> HR .
 \end{CD}
\end{equation}
Here $\rho_R$ inserts one ring coefficient into its selected
coordinate, and $\widetilde c$ does the same with $c$.  Their
naturality follows from the single-coordinate argument just
given; their multiplicativity and units follow from (HG19),
so (HG20) is proved coordinatewise.  This is a precise
multiplicative replacement, with its correction state visible.

Amplitude also has a universal characterization.  For any ring
$A$, every unital semiring homomorphism $u:G\to A$ factors
uniquely through $q:G\to R$.  Indeed every $z_\lambda$ is
additively idempotent, so $u(z_\lambda)=0$ in the additive
group of $A$.  Therefore $v(r)=u(\widehat r)$ satisfies
$v(0)=0$, $v(1)=1$, and preserves sums and products by (HG1),
giving a ring map $v:R\to A$ with $u=vq$.  More generally this
same proof, omitting multiplication, shows that $q$ is the
universal map from the additive monoid of $G$ into an abelian
group.  Thus
\begin{equation}\tag{HG21}
 G^{\rm gp}_{\rm additive}\cong(R,+),\qquad
 \operatorname{Hom}_{\rm semiring}(G,A)
       \cong\operatorname{Hom}_{\rm ring}(R,A).
\end{equation}
This explains exactly which coefficient information a ring-valued
realization forgets: the full ideal (HG9).  It does not provide
a $\Gamma$-section in the reverse direction.

\begin{proposition}\label{HG:completion}
Degreewise additive Grothendieck group completion gives the exact
natural isomorphism
\begin{equation}\tag{HG21a}
 (HG(X))^{\rm gp}\ \xrightarrow{\ \cong\ }\ HR(X),
 \qquad [(r,\lambda)]\longmapsto r.
\end{equation}
It preserves all pointed-set maps and the bilinear multiplication
after group completion.  The additive congruence that effects this
quotient is exactly
\begin{equation}\tag{HG21b}
 (r,\lambda)\sim(s,\mu)
 \quad\Longleftrightarrow\quad r=s
 \quad\Longleftrightarrow\quad
 (r,\lambda)+(e)_{x\in X^\circ}
  =(s,\mu)+(e)_{x\in X^\circ}.
\end{equation}
Thus group completion identifies all support choices over each
fixed amplitude tuple and makes no further identifications.
\end{proposition}
\begin{proof}
For a single coordinate, additive idempotence gives
$[z_\lambda]=[z_\lambda]+[z_\lambda]$, hence
$[z_\lambda]=0$ in a group.  The identity
$\widehat r+\widehat{-r}=e$ gives
$[\widehat{-r}]=-[\widehat r]$.  The map
$r\mapsto[\widehat r]$ is additive, because the sum of any
two top-supported coefficients is the top-supported coefficient
of their ring sum, including cancellation.  Its value at zero
is $[e]=0$.  It is inverse to the map induced by $q$.
For finitely many coordinates, every tuple is the sum of its
one-coordinate insertions; the same argument therefore gives
the inverse $r\mapsto\sum_x[\widehat{r_x}\text{ at }x,
\tau\text{ elsewhere}]$ and proves (HG21a) with its universal
property.  More explicitly, any additive map from $G^{X^\circ}$
to an abelian group kills every tuple of labelled zeros and is
determined uniquely by these coordinate classes, so factors
uniquely through $q^{X^\circ}$.

Adding $e$ at a coordinate replaces its support by $1_L$
without changing its amplitude.  This proves (HG21b).
Conversely the congruence must identify all labelled zeros
because their classes vanish; its remaining quotient is
already the group $R^{X^\circ}$ by the displayed inverse,
so no larger congruence is present.

Applying a pointed-set map to a tuple class replaces the
amplitude at $y$ by $\sum_{f(x)=y}r_x$ by (HG3), which is
exactly $HR(f)$.  The isomorphism is therefore natural on
the entire category of finite pointed sets, not only on
projections.  The product on the additive completions is
determined by its values on coefficient classes, where (HG4)
becomes $r_xs_y$; its bilinear extension is precisely the
usual product of $HR$.  This proves the asserted compatibility.
\end{proof}

The conclusion is the following pair of proved maps, with the
operation on the right explicitly stated:
\begin{equation}\tag{HG21c}
 HL\ \xrightarrow{\ \iota\ }\ HG_L(R)
 \ \xrightarrow{\ Hq\ }\ HR
       \cong \bigl(HG_L(R)\bigr)^{\rm gp}_{\rm degreewise}.
\end{equation}
It is not a claim that an arbitrary stable realization of
$HG_L(R)$ is equivalent to the Eilenberg--Mac Lane spectrum
of $R$.  It is also not an unmentioned group completion of
the bounded objects introduced next: those bounds need not
be closed under coordinatewise addition.

\subsection{The original seminorm constraint retains the full support ideal}

Give $R$ a submultiplicative seminorm with the exact axioms in
Connes--Consani (\texttt{subverti1}), source lines 895--903:
$|0|=0$, $|1|=1$, $|r+s|\le |r|+|s|$, and
$|rs|\le |r||s|$.  Define
\begin{equation}\tag{HG22}
 |(r,\lambda)|_G=|r|.
\end{equation}
This is a submultiplicative seminorm on $G$, as each required
inequality reduces to its amplitude counterpart.  For $c>0$
define the closed and strict functors
\begin{equation}\tag{HG23}
 \begin{aligned}
 B_c^L(X)&=\{(r,\lambda)\in HG(X):\sum_{x\in X^\circ}|r_x|\le c\},\\
 B_{<c}^L(X)&=\{(r,\lambda)\in HG(X):\sum_{x\in X^\circ}|r_x|<c\}.
 \end{aligned}
\end{equation}
\begin{theorem}\label{HG:bound}
$B_1^L$ is a unital $\mathbb S$-subalgebra of $HG$.
Both functors in (HG23) are modules over $B_1^L$ and the
multiplication maps restrict to
\begin{equation}\tag{HG24}
 B_c^L\wedge B_d^L\to B_{cd}^L,
 \qquad B_{<c}^L\wedge B_d^L\to B_{<cd}^L
 \quad(c,d>0).
\end{equation}
Under amplitude the fibres are still exactly (HG8), for every
amplitude tuple obeying its bound.  Under support the fibre at
$\lambda$ is exactly
\begin{equation}\tag{HG25}
 \{(r_x)_{x:\lambda_x=1_L}:
       \sum_{x:\lambda_x=1_L}|r_x|\le c\},
\end{equation}
with the corresponding strict inequality for $B_{<c}^L$.
The entire $HL$ embeds into every one of these functors, and
is precisely the zero-amplitude fibre.  In particular every
support label has zero seminorm cost.
\end{theorem}
\begin{proof}
The fibre-sum inequality in the source becomes, with all supports
retained,
\begin{equation}\tag{HG26}
 \sum_{y\in Y^\circ}\left|\sum_{x:f(x)=y}r_x\right|
 \le\sum_{x:f(x)\ne *}|r_x|
 \le\sum_{x\in X^\circ}|r_x|.
\end{equation}
Thus (HG3) preserves both bounds.  The exact support joins
remain those of (HG3), since no inequality has changed any
coefficient.  Likewise
\begin{equation}\tag{HG27}
 \sum_{x,y}|r_xs_y|\le
         \left(\sum_x|r_x|\right)\left(\sum_y|s_y|\right)
\end{equation}
proves (HG24), including the strict case since $d>0$.
Unit tuples have seminorm sum $1$, so are in $B_1^L$.
The module identities are the restrictions of the already
proved algebra identities.  Since (HG23) depends only on
amplitude, it removes none of the support choices over any
allowed tuple.  Substituting the support fibre from (HG8)
gives (HG25).  Zero tuples have total seminorm zero, so all
of $HL$ remains in every positive closed or strict bound.
\end{proof}

Let $N=\{r\in R:|r|=0\}$.  The triangle and product
inequalities make $N$ an ideal: closure under additive
inverses follows from $|-r|\le|-1||r|=0$.  At the closed
zero bound the exact object is
\begin{equation}\tag{HG28}
 B_0^L(X)=\{(r,\lambda)\in HG(X):r_x\in N\text{ for all }x\}.
\end{equation}
It equals $HL$ precisely when $N=\{0\}$.  This distinction
prevents a degenerate seminorm from silently erasing nonzero
amplitudes.  The strict unit ball is not unital, because the
unit tuple has seminorm sum exactly $1$; it still retains
every element of $HL$.

\begin{theorem}\label{HG:boundedsection}
Assume $0_L\ne1_L$ and let $B_c(R)$ denote the source's
classical amplitude functor with the same bound $c>0$.
A $\Gamma$-section of $q:B_c^L\to B_c(R)$ exists exactly
when
\begin{equation}\tag{HG34}
 |r|+|-r|>c\qquad\text{for every nonzero }r\in R.
\end{equation}
When it exists it is unique and is the coordinatewise map
$b_X$.  For strict balls $B_{<c}$ the exact criterion is
$|r|+|-r|\ge c$ for all nonzero $r$.  At $c=1$, an existing
$\Gamma$-section is a unital $\mathbb S$-algebra section
exactly when, in addition,
\begin{equation}\tag{HG35}
 r,s\ne0,\quad |r|\le1,\quad |s|\le1
             \quad\Longrightarrow\quad rs\ne0.
\end{equation}
\end{theorem}
\begin{proof}
Projection to any coordinate is a map of the bounded functors,
because discarding terms cannot increase the sum of seminorms.
Thus any section is determined coordinatewise by its value
at $1_+$.  A nonzero amplitude has only one preimage in $G$,
and pointedness forces the zero amplitude to go to $\tau$.
Therefore $b_X$ is the only possible section.

If $r\ne0$ and $|r|+|-r|\le c$, the tuple $(r,-r)$ is
admissible; its two paths in the figure following (HG14) have
endpoints $e$ and $\tau$.  Naturality fails.  Conversely,
suppose (HG34).  A nonzero correction $C_f(a)_y$ for an
admissible tuple would give a zero-sum fibre with some nonzero
summand $a_i$.  In the ring,
$-a_i=\sum_{x\in I_y\setminus\{i\}}a_x$.  The triangle
inequality gives
\[
 |a_i|+|-a_i|\le\sum_{x\in I_y}|a_x|
        \le\sum_{x\in X^\circ}|a_x|\le c,
\]
contrary to (HG34).  Hence every correction vanishes and
(HG14) proves naturality.  The strict case has the same
proof with the final strict inequality; its failed pair
condition is correspondingly $|r|+|-r|<c$.

Unit preservation holds for the coordinatewise $b$, and
(HG16) shows that multiplication is preserved if (HG35)
holds, because all individual amplitudes in a unit-ball
tuple have seminorm at most $1$.  If (HG35) fails, place
the two nonzero zero divisors at level $1$.  Each is allowed
by the unit bound and their product produces $e$ instead
of $\tau$, proving failure of multiplication.  This proves
both directions of the criterion.
\end{proof}

\subsection{The sheaf calculation on the source's compactified base}

The indexing space here is the source's
$X=\overline{\operatorname{Spec}\mathbb Z}$.  Its topology is
defined in \texttt{sectstructsh}, lines 938--946.  This section
constructs a coefficient sheaf on that specified space; no
identification with a prime spectrum of $G_L(\mathbb Z)$ is used.
For a nonempty open $U$ put
\begin{equation}\tag{HG29}
 R_U=\{r\in\mathbb Q:v_p(r)\ge0\text{ for every finite prime }p\in U\}.
\end{equation}
Define a sheaf candidate, with empty-open value the one-point
algebra, by
\begin{equation}\tag{HG30}
 \mathcal A_L(U)(K)=
 \begin{cases}
 HG_L(R_U)(K),&\infty\notin U,\\
 \{(r,\lambda)\in HG_L(R_U)(K):\sum_{k\in K^\circ}|r_k|\le1\},
       &\infty\in U.
 \end{cases}
\end{equation}
The absolute value is the usual one on $\mathbb Q$.  The restriction
map for nonempty $V\subset U$ includes $R_U$ into $R_V$ and
retains each support coordinate unchanged.

\begin{theorem}\label{HG:sheaf}
Equation (HG30) is a sheaf of unital $\mathbb S$-algebras, with
amplitude morphism to the exact Connes--Consani sheaf
$\mathcal O_X$ of (\texttt{subsalg}) and Proposition
\texttt{propspecz}, lines 944--977.  Its amplitude-zero sheaf
is the constant $HL$ on nonempty opens, including opens
containing $\infty$.  Its stalks are
\begin{equation}\tag{HG31}
 \begin{aligned}
 (\mathcal A_L)_\eta&=HG_L(\mathbb Q),\\
 (\mathcal A_L)_p&=HG_L(\mathbb Z_{(p)})\quad(p\text{ finite}),\\
 (\mathcal A_L)_\infty&=B_1^L(\mathbb Q).
 \end{aligned}
\end{equation}
All these identities retain the full lattice $L$.
\end{theorem}
\begin{proof}
Every value in (HG30) is a unital algebra by
Theorem~\ref{HG:functor} or Theorem~\ref{HG:bound}; the restrictions
are the stated algebra inclusions.  Every pair of nonempty
opens intersects in the generic point.  Thus compatible local
sections over a cover have literally the same rational
amplitudes and the same support coordinates in the ambient
$HG_L(\mathbb Q)(K)$.  An amplitude integral at all primes
in each covering open is integral at every prime of their
union.  If the union contains $\infty$, one covering open
does too, and imposes the norm bound on that common tuple.
The common tuple therefore belongs to (HG30) on the union,
and is its unique gluing.  This proves the sheaf condition at
every $K$, with multiplication and units already compatible.
Amplitude removes precisely the support coordinates and gives
the source sheaf formulas.  Its zero fibre is $HL$ on every
nonempty open, with identity restriction maps.

At the generic point one can take opens excluding $\infty$
and any finite set of denominator primes, so every rational
tuple occurs.  At a finite prime $p$ the same procedure allows
all denominator primes other than $p$ and excludes $\infty$,
giving exactly $\mathbb Z_{(p)}$.  At $\infty$ every open
retains the norm bound, but any finite collection of rational
denominators is allowed after removing its finite prime
divisors.  Hence the union of these bounded sets is exactly
$B_1^L(\mathbb Q)(K)$ at each $K$.  The restrictions never
change supports, so the stalk descriptions include all of $L$.
\end{proof}

The global sections at a finite pointed set $k_+$ are especially
explicit.  Since $R_X=\mathbb Z$, the inequality
$\sum_{i=1}^k|r_i|\le1$ allows either the zero amplitude
tuple or a single coordinate with value $+1$ or $-1$.
Consequently
\begin{equation}\tag{HG32}
 \mathcal A_L(X)(k_+)=
 L^k\ \sqcup\
 \coprod_{i=1}^k\coprod_{\varepsilon\in\{+1,-1\}}
       \bigl\{r_i=\varepsilon,\ \lambda_i=1_L,\
              r_j=0,\ \lambda_j\in L\ (j\ne i)\bigr\}.
\end{equation}
These are disjoint subsets; the first is the entire
zero-amplitude fibre, whose own pointed zero is the all-bottom
tuple.  If $L$ is finite of cardinality $N$, then for $k\ge1$
\begin{equation}\tag{HG33}
 |\mathcal A_L(X)(k_+)|=N^k+2kN^{k-1}.
\end{equation}
For $k=0$ there is one section.  At level $1$, the global
sections are all $z_\lambda$ together with $\widehat1$ and
$\widehat{-1}$, so their number is $N+2$ when $L$ is finite.
This is a direct application of the original unit bound, with
no added positivity assumption.

\begin{proposition}\label{HG:globalstalkdefect}
For $0_L\ne1_L$, the global-sections amplitude map
\begin{equation}\tag{HG36}
 \mathcal A_L(X)\longrightarrow\mathcal O_X(X)
\end{equation}
has a unique unital $\mathbb S$-algebra section, given by
$b$ in every coordinate.  Nevertheless the sheaf amplitude
map $\mathcal A_L\to\mathcal O_X$ has no section even as
a sheaf of $\Gamma$-sets.  Already on the single open
$U=X\setminus\{2\}$ there is no $\Gamma$-section.  There
is likewise no $\Gamma$-section at the infinity stalk.
\end{proposition}
\begin{proof}
For a nonzero integer $r$, $|r|+|-r|=2|r|\ge2>1$,
and $\mathbb Z$ has no nonzero zero divisors.  Theorems
\ref{HG:bound} and~\ref{HG:boundedsection} therefore prove
the unique algebra section on global sections.  Its image
has all unused coordinates equal to $\tau$, leaving the
other labelled-zero points outside this chosen section.

On $U=X\setminus\{2\}$ the coefficient ring is
$\mathbb Z[1/2]$ and the archimedean bound is still exactly
$1$.  Thus $(1/2,-1/2)$ is an allowed tuple and
\begin{equation}\tag{HG37}
 HG(\mathrm{fold})(\widehat{1/2},\widehat{-1/2})=e,
 \qquad b(HR(\mathrm{fold})(1/2,-1/2))=\tau.
\end{equation}
Pointedness and the coordinate projections force any
putative section to use those two lifts, so there is no
alternative choice that could repair this failure.  A
section of sheaves would evaluate to a section on $U$,
which is impossible.  The same allowed tuple in
$B_1(\mathbb Q)$, the source infinity stalk from (HG31),
proves its asserted failure directly.
\end{proof}

For an arbitrary bound $c>0$, the same proof replaces (HG32)
by all integer vectors satisfying $\sum_i|r_i|\le c$, each
with fibre $L^{\{i:r_i=0\}}$.  For $0<c<1$ the entire global
bounded object is exactly $HL$; at the strict bound $c=1$ it
is again exactly $HL$.  At a rational or complex archimedean
coefficient ring, a pair $r,-r$ with $r\ne0$ and $2|r|\le1$
is allowed by the unit ball.  Folding its two top-supported
coordinates gives $e$, whereas folding the zero tuple gives
$\tau$.  Thus cancellation-labelled zero survives the actual
archimedean restriction at precisely the original norm bound.

The complete retained object is therefore (HG7) with (HG3),
(HG4), and (HG23), and its nontrivial zero-amplitude ideal is
(HG9)--(HG10).  The original seminorm does not force that ideal
to vanish.  Ring-valued realization removes it by (HG21),
while the $\Gamma$-algebra and every positive archimedean bound
retain it explicitly.

\subsection{Exact comparison with the later all-subset Riemann--Roch module}

The following second source uses a different, explicitly stated
module: Alain Connes and Caterina Consani,
\emph{Riemann--Roch for the ring $\mathbb Z$},
\href{https://arxiv.org/abs/2306.00456}{arXiv:2306.00456},
original \nolinkurl{RRZ.tex}, equation \texttt{normdefn},
Lemma \texttt{xlem}, and equation \texttt{xdefn}, source
lines 185--208.  That source records the older sum-of-absolute-values
bound and then defines the new module by all subset sums.
These definitions give an exact inclusion without identifying
the two modules or changing either original bound.

For an integer $m\ge0$ set
\begin{equation}\tag{HG38}
 \begin{aligned}
 E_{L,m}(k_+)=\{((a_i,\lambda_i))\in HG_L(\mathbb Z)(k_+):
     &-m\le\sum_{i\in I}a_i\le m\\
     &\text{for every }I\subseteq\{1,\ldots,k\}\}.
 \end{aligned}
\end{equation}
This is the same complete support lift as (SRR2) in the
accompanying programme calculation, here defined in full so
that the comparison needs none of its dimension claims.
An output subset under a pointed map has as preimage a subset
of input coordinates, so (HG3) preserves (HG38).  Thus it
is a $\Gamma$-subfunctor.  Its amplitude version is denoted
$E_m$.

\begin{proposition}\label{HG:twobounds}
For $P(a)=\sum_{a_i>0}a_i$ and
$N(a)=\sum_{a_i<0}(-a_i)$, the two exact bounds are
\begin{equation}\tag{HG39}
 \begin{aligned}
 (a,\lambda)\in B_m^L(\mathbb Z)(k_+)
       &\ \Longleftrightarrow\ P(a)+N(a)\le m,\\
 (a,\lambda)\in E_{L,m}(k_+)
       &\ \Longleftrightarrow\ P(a)\le m\text{ and }N(a)\le m.
 \end{aligned}
\end{equation}
There are natural inclusions preserving every amplitude and
support coordinate:
\begin{equation}\tag{HG40}
 B_m^L(\mathbb Z)\ \lhook\joinrel\longrightarrow\ E_{L,m}
       \ \lhook\joinrel\longrightarrow\ B_{2m}^L(\mathbb Z).
\end{equation}
For $m\ge1$ both inclusions are strict; for $m=0$ all three
objects are exactly $HL$.  The amplitude diagrams for these
inclusions commute, and their zero-amplitude fibres are
identically $HL$, with no quotient of the lattice labels.
\end{proposition}
\begin{proof}
The first equivalence is the identity
$\sum_i|a_i|=P(a)+N(a)$.  Every subset sum is between
$-N(a)$ and $P(a)$, and these endpoints are attained by
taking all negative or all positive coordinates respectively.
This proves the second equivalence, including empty subsets,
and is the calculation in the source's subsequent
\texttt{xdefnr}.  The numerical implications
$P+N\le m\Rightarrow(P\le m,N\le m)\Rightarrow P+N\le2m$
prove (HG40).  All three use exactly the same support
constraint (HG2) and the same pushforward (HG3), proving
naturality and compatibility with amplitude.

The tuple $(\widehat m,\widehat{-m})$ belongs to $E_{L,m}$
but has absolute-value sum $2m>m$, so is outside the first
ball.  The one-coordinate tuple $\widehat{2m}$ belongs to
the second ball but its singleton subset sum exceeds $m$,
so is outside $E_{L,m}$.  At $m=0$, each individual
coordinate is a subset sum, hence every amplitude is zero;
all support tuples remain, giving $HL$.  The same observation
proves the assertion about zero-amplitude fibres for every
$m$.
\end{proof}

\begin{proposition}\label{HG:rrsectiondefect}
For $0_L\ne1_L$ and every $m\ge1$, the map
$E_{L,m}\to E_m$ has no $\Gamma$-section.  In particular,
the unique global $\Gamma$-algebra section on the older
$B_1(\mathbb Z)$ from (HG36) does not extend to the later
$E_1$ module, even though their level-$1$ point sets agree.
\end{proposition}
\begin{proof}
Pointed coordinate projections are available in the all-subset
module.  Therefore any section must send zero to $\tau$
and every nonzero coordinate $a$ to its unique lift
$\widehat a$, exactly as in the proof of
Theorem~\ref{HG:boundedsection}.  For every $m\ge1$, the
tuple $(1,-1)$ belongs to $E_m(2_+)$: its subset sums are
$0,1,-1,0$, all in $[-m,m]$.  The required naturality
equation would be
\begin{equation}\tag{HG41}
 e=HG(\mathrm{fold})(\widehat1,\widehat{-1})
       =b(E_m(\mathrm{fold})(1,-1))=b(0)=\tau,
\end{equation}
which is false for $0_L\ne1_L$.  At $m=1$ the tuple
$(1,-1)$ is exactly an additional tuple supplied by the
first strict inclusion in (HG40); it was disallowed by
the older absolute-value sum bound.  At level $1$ both
bounds require only $|a|\le1$, so their point sets do
agree there.  The failure therefore occurs in their
actual higher-level fold structures, as calculated, and
does not change any coordinate or bound.
\end{proof}
